\begin{lemma}
\label{lemma-functorial-homotopies}
Let $\mathcal{A}$ be a Grothendieck abelian category.
Let $(K_i^\bullet)_{i \in I}$ be a set of acyclic complexes.
There exists a functor $M^\bullet \mapsto \mathbf{M}^\bullet(M^\bullet)$
and a natural transformation
$j_{M^\bullet} : M^\bullet \to \mathbf{M}^\bullet(M^\bullet)$
such
\begin{enumerate}
\item $j_{M^\bullet}$ is a (termwise) injective quasi-isomorphism, and
\item for every $i \in I$ and $w : K_i^\bullet \to M^\bullet$
the morphism $j_{M^\bullet} \circ w$ is homotopic to zero.
\end{enumerate}
\end{lemma}

\begin{proof}
For every $i \in I$ choose a (termwise) injective map of complexes
$K_i^\bullet \to L_i^\bullet$ which is homotopic to zero with
$L_i^\bullet$ quasi-isomorphic to zero. For example, take $L_i^\bullet$
to be the cone on the identity of $K_i^\bullet$.
We define $\mathbf{M}^\bullet(M^\bullet)$ by the following pushout diagram
$$
\xymatrix{
\bigoplus_{i \in I}
\bigoplus_{w : K_i^\bullet \to M^\bullet}
K_i^\bullet \ar[r] \ar[d] & M^\bullet \ar[d] \\
\bigoplus_{i \in I}
\bigoplus_{w : K_i^\bullet \to M^\bullet}
L_i^\bullet \ar[r] &  \mathbf{M}^\bullet(M^\bullet).
}
$$
Then $M^\bullet \mapsto \mathbf{M}^\bullet(M^\bullet)$ is a functor. The right
vertical arrow defines the functorial injective map $j_{M^\bullet}$.
The cokernel of $j_{M^\bullet}$ is isomorphic to the direct sum of
the cokernels of the maps $K_i^\bullet \to L_i^\bullet$ hence acyclic.
Thus $j_{M^\bullet}$ is a quasi-isomorphism. Part (2) holds by construction.
\end{proof}
